\section{Discussion}
The simulation results are not surprising in isolation (clustered data need clustered inference), but their size is. Reporting a single pooled interval from several runs is common and, with ten runs per task, wrong by almost thirty coverage points. The public re-analysis shows that the consequence for leaderboards is concrete: with 500 instances and one run each, differences of up to about five points are routinely within noise, and the neighbours on the leaderboard are not separated. Neither finding says that leaderboards are useless: large gaps are reliably detected, and the top of the SWE-bench ranking is clearly above the bottom. They say that rank positions near each other should not be read as an ordering.

\paragraph{Recommendations.} (i) Treat tasks as the sampling unit; report task-bootstrap or hierarchical intervals. (ii) Compare agents on the same tasks with paired task-level tests and correct for multiplicity. (iii) Plan with the variance decomposition: more tasks before more repeats, unless reliability (pass$^k$) is the target or tasks are expensive relative to runs. (iv) Report pass$^k$ with pass@$k$. (v) If monitoring an evaluation as it runs, use a confidence sequence, not repeated fixed-sample tests. (vi) Publish per-task, per-run outcomes so that others can do this.

\paragraph{Refuted or weakened parts of the hypothesis.} \REALREFUTED
\section{Limitations, Ethics and Broader Impact}\label{sec:limits}
\paragraph{Limitations.} The local-model evaluation uses small quantised models on one machine, so absolute scores and failure modes do not transfer to frontier agents; generality rests on the simulations and the public re-analysis. The environments are cleaner than real software tasks and checkers verify only the state the generator anticipates. Tasks are i.i.d.\ draws from the population defined by our generator, so intervals refer to that population, not to ``ability''. The variance decomposition is linear on the probability scale. The SWE-bench re-analysis assumes full coverage of the 500 instances by every submission and ignores differences in scaffolds, budgets and submission dates; the tau-bench data cover two models. Confidence sequences are conservative. Task difficulty in the benchmark is controlled by knobs we chose and is not calibrated to any human notion of difficulty.
\paragraph{Ethics and impact.} No human subjects or personal data are involved; all environments are synthetic and all public data are used under their licenses (tau-bench: MIT; the SWE-bench experiments repository declares no license, so we store only derived aggregates and a download script). Better error bars can reduce overclaiming in agent evaluation, but statistical rigour can also be used to dismiss real improvements; a non-significant result is not evidence of equivalence, and we provide equivalence-margin stopping for that reason. All experiments used free, local, open-weight models and public data; no paid API was used.
\section{Conclusion}
Agent leaderboards carry less information than their precision suggests. A statistically principled treatment of tasks, runs and their interaction, tested against ground truth, shows where differences are real and tells practitioners how to spend an evaluation budget. We release the library, the benchmark, the playground and all scripts to regenerate every number in this paper.
